\subsection{Bounded-weight free modules and an exact flatness calculation}

\begin{proposition}[A-L14: Proved; bounded-below free graded modules]
\label{AL14}
Let $p$ be a prime, $D=\Gamma_{\mathbb F_p}(T)$, and
\[
 F=\bigoplus_{\lambda\in L}D\langle a_\lambda\rangle,
 \qquad a_\lambda\geq A
 \tag{AL.78}
\]
for one integer $A$. The generator of $D\langle a\rangle$
has weight $a$. There is no restriction on $L$ or on the
number of generators in a weight. Then $F$ is injective
in the category of integer-weight-graded $D$-modules
and degree-zero homomorphisms.
\end{proposition}

\begin{proof}
We use Nagpal--Snowden, Theorem~7.1(a), statement on
printed page~39 and proof on page~41 of the opened v2
source \cite{IC-nagpal-snowden}: $D$ itself is graded
injective. The classical divided power parameters vanish
at every positive power of $p$, so the theorem applies.
The unqualified uniqueness sentence in that proof fails
for a zero restriction. Indeed for $I=(T)\subset D$
the zero map $I\langle p-1\rangle\to D$ has distinct
extensions $D\langle p-1\rangle\to D$ given by
$1\mapsto0$ and $1\mapsto T^{[p-1]}$:
$T T^{[p-1]}=pT^{[p]}=0$, while $T^{[p-1]}\ne0$.
The theorem's extension-existence conclusion is obtained
by the following precise qualification. The uniqueness
argument is used on a nonzero restriction: a nonzero map from a
homogeneous ideal has a nonzero restriction to some
finitely generated subideal. Extensions of that restriction
are unique because the relevant homogeneous component
of $D$ is at most one-dimensional and the restriction
map is nonzero. Every larger finitely generated subideal
therefore gives the same extension. A zero ideal map has
the zero extension. This supplies exactly the extension
property of that theorem used here.

Let $I\subset D$ be a homogeneous ideal and let
$f:I\langle d\rangle\to F$ be a degree-zero map.
For each $\lambda$, extend its coordinate map to
$D\langle d\rangle\to D\langle a_\lambda\rangle$.
Choose the zero extension when the coordinate map is
zero. The image of the generator in a nonzero extension
is
\[
 c_\lambda T^{[k_\lambda]},\qquad
 k_\lambda=d-a_\lambda\in\{0,\ldots,d-A\},
 \quad c_\lambda\in\mathbb F_p^\times .
 \tag{AL.79}
\]
An empty interval means that $f=0$.
For each integer $k$ in this finite interval with
$I T^{[k]}\ne0$, choose one homogeneous $z_k\in I$
with $z_kT^{[k]}\ne0$. Every nonzero coordinate map
whose extension has exponent $k$ is nonzero on $z_k$.
It therefore belongs to the finite coordinate support
of $f(z_k)\in F$. The union of these supports over the
finitely many $k$ is finite. Thus only finitely many
coordinate extensions are nonzero; their sum is an
actual map $D\langle d\rangle\to F$ extending $f$.

This ideal-extension property gives injectivity directly.
For an inclusion of graded modules $M\subset N$ and a
map $M\to F$, take a maximal graded extension by Zorn's
lemma. A homogeneous element $x$ outside its domain has
a homogeneous ideal of coefficients $r$ with $rx$ in
that domain. The already defined values of $rx$ give
an ideal map of the preceding kind; an extension of
that map defines the image of $x$, and hence an extension
to the submodule generated by the domain and $x$.
Maximality forces the domain to be $N$.

The explicit direct sum in Nagpal--Snowden, Remark~7.6,
page~41, has generator weights $1-n_i$ for an unbounded
sequence $n_i$. Its weights are unbounded below. The
finite interval in (AL.79) states the exact restriction
used in the present proof.
\end{proof}

\begin{proposition}[A-L15: Proved; low boundaries over the integral coefficient algebra]
\label{AL15}
Let $R$ be $\mathbb Z$ or $\mathbb Z_{(p)}$, and put
$S=\Gamma_R(T)$. Let $M^\bullet$ be a nonnegative
cochain complex of weight-graded $S$-modules. Each $M^i$
is graded free with generator weights bounded below;
the bound may depend on $i$. Write
$Z^i=\ker(d:M^i\to M^{i+1})$ and
$B^i=\operatorname{im}(d:M^{i-1}\to M^i)$.
Suppose the given coefficient inclusion
$\iota:S\to M^0$ has image $Z^0$, and comes with an
$S$-linear graded retraction $e:M^0\to S$ satisfying
$e\iota=\mathrm{id}$. Let $d_0\geq0$ and suppose
$H^i(M)=0$ for $1\leq i\leq d_0$.
Then all of
\[
 B^1,\ldots,B^{d_0+1},\qquad
 Z^0,Z^1,\ldots,Z^{d_0}
 \tag{AL.80}
\]
are flat as ungraded $S$-modules. For every residue
prime $\ell$ of $R$, each $B^i/\ell B^i$ in this
range is a graded direct summand of a graded free
$\Gamma_{\mathbb F_\ell}(T)$-module whose generator
weights are bounded below. No assertion about
$Z^{d_0+1}$ is included.
\end{proposition}

\begin{proof}
First take $R=\mathbb Z_{(p)}$. Every $M^i$, $Z^i$
and $B^i$ is $R$-torsionfree; the last two are actual
submodules of free $R$-modules. The split inclusion in
degree zero identifies $B^1=M^0/\iota(S)$ with the
graded direct summand $\ker e$ of $M^0$. In particular
$B^1/pB^1$ is a graded projective and, by A-L14 and
the splitting, a graded injective module over $S/pS$.

For $1\leq i\leq d_0$ the cohomology hypothesis gives
$B^i=Z^i$, and the actual sequence
\[
 0\longrightarrow B^i\longrightarrow M^i
       \xrightarrow{d} B^{i+1}\longrightarrow0
 \tag{AL.81}
\]
is exact. Its quotient is $R$-torsionfree, so reduction
modulo $p$ preserves this exact sequence. Inductively,
graded injectivity of $B^i/pB^i$ splits the reduced
inclusion. Hence $B^{i+1}/pB^{i+1}$ is a graded direct
summand of $M^i/pM^i$. A-L14 makes this new summand
graded injective as well as projective, continuing the
induction. After forgetting weights it remains a direct
summand of an ordinary free module, and is therefore flat.

Rationally, every $B^i_{\mathbb Q}$ is a submodule of
the free $\mathbb Q[T]$-module $M^i_{\mathbb Q}$.
Every finitely generated submodule of it is torsionfree
over the principal ideal domain $\mathbb Q[T]$, hence
free. The union over these submodules is filtered;
exactness of filtered unions shows that $B^i_{\mathbb Q}$
is flat. We now supply the integral passage completely.

Let $L$ be any $S$-module that is $R$-torsionfree,
with $L/pL$ flat over $S/pS$ and $L_{\mathbb Q}$ flat
over $S_{\mathbb Q}$. Take a free $S$-resolution
$P_\bullet\twoheadrightarrow L$.
The length-one free resolution
$0\to S\xrightarrow{p}S\to S/pS\to0$ shows that
$\operatorname{Tor}^S_j(L,S/pS)=0$ for $j>0$:
the only possible positive group is the kernel of $p$
on $L$, which vanishes. Consequently
$P_\bullet/pP_\bullet$ is a free $S/pS$-resolution of
$L/pL$. For every $S$-module $N$ killed by $p$, the
actual equality of tensor complexes
\[
 P_\bullet\otimes_S N
 = (P_\bullet/pP_\bullet)\otimes_{S/pS}N
 \tag{AL.82}
\]
therefore gives $\operatorname{Tor}^S_j(L,N)=0$ for
every $j>0$. A finite filtration gives the same result
for modules killed by $p^r$. Taking the filtered union
of their $p^r$-torsion submodules gives the result for
every $p$-power-torsion module. This passage commutes
with homology of $P_\bullet\otimes_S-$ because filtered
colimits of modules are exact.

For a $\mathbb Q$-module $N$, localization gives the
actual tensor complex
$P_{\bullet,\mathbb Q}\otimes_{S_{\mathbb Q}}N$;
flatness of $L_{\mathbb Q}$ again gives vanishing of
all its positive homology. For an $R$-torsionfree
module $N$, the exact sequence
\[
 0\longrightarrow N\longrightarrow N_{\mathbb Q}
       \longrightarrow N_{\mathbb Q}/N\longrightarrow0
 \tag{AL.83}
\]
has $p$-power-torsion quotient. The Tor exact sequence
and the two vanishing results just proved give
$\operatorname{Tor}^S_j(L,N)=0$ for every $j>0$.
Finally every $N$ is an extension of its
$p$-power-torsion submodule by an $R$-torsionfree
quotient, so the same vanishing holds for arbitrary
$N$. Vanishing of $\operatorname{Tor}^S_1(L,-)$ is
the flatness assertion: applying the tensor--Tor exact
sequence to any short exact sequence proves exactness
of $L\otimes_S-$.

Apply this argument to each $L=B^i$ obtained above.
It proves (AL.80) for $R=\mathbb Z_{(p)}$, including
$d_0=0$. For $R=\mathbb Z$, localize the entire given
complex, splitting and exact sequences at each rational
prime. The preceding proof applies to every localization.
For arbitrary $N$, the abelian group
$\operatorname{Tor}^S_1(B^i,N)$ consequently localizes
to zero at every rational prime. An abelian group with
that property is zero: a nonzero element has either
infinite order, or finite order divisible by a prime
at whose localization it survives. Thus $B^i$ is
$S$-flat. The equality $Z^i=B^i$ in the stated range,
and $Z^0=\iota(S)$, prove the remaining assertions.
\end{proof}

These propositions use actual freeness of the coefficient
modules and actual global cohomological exactness in the
listed degrees. A-S09-24 supplies the latter finite-degree
condition for the reduced main gate. An $S$-free replacement
with actual comparison maps for every object of that gate
has not been constructed here.
